\chapter{Infrared Limits That Follow from the Proved Chain}
\label{app:gr_recovery}

This appendix collects the infrared limits that follow directly from
the theorems and structural results already established in
Chapters~3--12. It does not treat open problems as closed, and it
does not claim additional recovery of general relativity beyond
those limits.

\section{Limits that are obtained}
\label{sec:recovery-obtained}

Under the named assumptions of the monograph the following hold.

\begin{enumerate}
  \item \textbf{High-acceleration Einstein equation.}
        When the reversible entanglement stress is negligible
        compared with ordinary matter stress, the infrared equation
        reduces to
        \begin{equation}
          G_{\mu\nu}+\Lambda_{\mathrm{eff}}g_{\mu\nu}
          =8\pi G_{\mathrm{eff}}T_{\mu\nu}^{\mathrm{matter}}.
        \end{equation}
        This is the content of the high-acceleration limit of
        the equation already written after Theorem~\ref{thm:conservation}.

  \item \textbf{Massless tensor polarizations.}
        The entanglement-equilibrium condition of Chapters~6--7
        forces \(m_T=0\), so the infrared spectrum contains the
        two transverse-traceless modes of a massless spin-2 field.
        This is Theorem~\ref{thm:mT}; it is not re-derived here.

  \item \textbf{Conservation.}
        The total stress (matter + tensor + entanglement) is
        divergenceless,
        \begin{equation}
          \nabla^\mu\bigl(T_{\mu\nu}^{\mathrm{matter}}
          +T_{\mu\nu}^{\mathrm{tensor}}
          +T_{\mu\nu}^{\mathrm{ent}}\bigr)=0,
        \end{equation}
        by Theorem~\ref{thm:conservation}. This identity constrains the
        total stress; it does not by itself derive any reduced drag used in
        later chapters, whose microscopic origin remains Problem~I.

  \item \textbf{Structural form of \(G_{\mathrm{eff}}\).}
        The conditional, truncated heat-kernel identity is recorded
        in Theorem~\ref{thm:Geff}; the physical coefficient
        calculation remains open in Problem~F.

  \item \textbf{Regulation of the quartic term.}
        The holographic boundary-entropy vacuum of Chapter~12
        supplies a consistent mechanism that can cancel the
        leading \(\Lambda_{\mathrm{UV}}^4\) contribution against
        a boundary term. The specific working value
        \(\Lambda_{\mathrm{eff}}=3/L_{\mathrm{IR}}^2\) is a choice
        compatible with that mechanism, not a derived necessity.
\end{enumerate}

\section{Limits that are not obtained}
\label{sec:recovery-not}

The following are frequently listed as recoveries but do not follow
from the proved chain.

\begin{itemize}
  \item Entanglement equilibrium from the microscopic path integral
        (Problem~B). The fixed-algebra first law does not derive
        \(\delta S_D=0\) for the non-conformal tensor block.

  \item Matrix-derived local Lorentz symmetry and a full field-content
        account, including the independent connection and ghost
        operators (Problem~D). Imposing local Lorentz symmetry gives
        only a scope-limited Presentation-B model.

  \item The infrared boundary entropy density and stationarity derived
        from the matrix theory (Problem~G). The conditional mechanism
        and an assumed boundary-model scaling do not derive these inputs.

  \item Exact \(\gamma-1=0\) in the solar system.
        The conformal mode \(\psi\) is not identical to the radial
        Higgs; its mass matrix has not been diagonalized
        (Problem~C). Cassini therefore remains a constraint on any
        light eigenmode that may survive.

  \item A first-principles drag exponent or numerical rate.
        The static conformal tail is conditional on its stated
        assumptions and does not determine a drag law. The assumed
        \(d^{-4}\) drag is a phenomenological input of Chapter~14
        (Problem~I). Its microscopic origin and compatibility with a
        covariantly conserved stress remain unresolved.

  \item Singularity resolution from the matrix theory.
        The regular-core profile is a phenomenological model. For
        \(\Lambda_{\mathrm{eff}}=0\) and \(Mr_0\ne0\) it is not an
        exact vacuum solution of the displayed first-integral equation
        in Appendix~\ref{app:status} (Problem~J), conditional on the
        stated reduction. That equation also excludes a nonzero
        integration constant for a smooth regular centre at any
        \(\Lambda_{\mathrm{eff}}\). Neither result derives the profile
        from the microscopic action.

  \item Lorentzian signature from the matrix algebra
        (Problem~E). It remains under the assumption stated in
        Chapter~3.

  \item An interpolating function for galactic or binary
        kinematics derived from \(I(A:B)\) (Problem~H).
        Coefficient-level compatibility with the published RAR
        map is recorded in Appendix~\ref{app:empirical}; it is
        not a derivation.
\end{itemize}

\section{Status}
\label{sec:recovery-status}

The infrared theory obtained from the proved chain is Einstein
gravity with a structural \(G_{\mathrm{eff}}\), a massless spin-2
sector, a divergenceless entanglement stress, and a consistent
regulation of the leading vacuum-energy term. Every stronger
recovery claim requires the closure of at least one open problem
listed in Appendix~\ref{app:status}; a completed matrix-derived
infrared derivation requires all the critical-path problems listed in
Chapter~\ref{ch:agenda} to be DERIVED, as specified there. This appendix adds no further
recovery.